\section{Conclusion}
\label{sec:conclusion}

We investigated whether a reinforcement-learning policy trained inside
RLTabucol learns guidance that is specific to its source framework or remains
useful when transferred to a structurally different search process. The
source-framework experiments show that the learned policy provides useful
vertex--color guidance relative to random recoloring, but RLTabucol does not
consistently outperform Pure Tabucol. The actor-only diagnostic further shows
that the policy can substantially reduce conflicts while Tabucol supplies most
of the fine repair, suggesting that the remaining limitation lies partly in how
the learned guidance is handed to the downstream solver rather than in the
absence of useful policy behavior.

Motivated by these results, we transferred the same policy to ACO
through a soft pheromone-injection interface. The transferred configuration
consistently improves on both unguided ACO and the random-policy transfer
reference at fixed $k$, and shows a stronger overall profile than the structured
conflict-based pheromone prior. The phase analysis further indicates that policy
utility is not limited to conflicts removed directly during the RL phase: the
learned color preferences can also guide the subsequent ACO construction
process. Soft transfer therefore preserves useful policy information while
allowing multiple policy-informed constructions instead of relying on a single
policy-produced handoff state.

The main optimization result is obtained by
RLACO$\rightarrow$Tabucol. The best coloring produced by the complete RLACO
stage is used to initialize a separate downstream Tabucol solver. This
configuration reaches the best or tied-best $k_{\mathrm{best}}$ on all eleven
tested instances, including two cases in which it reaches a lower color budget
than all three source-framework comparators. It also recovers the three cases
in which RLTabucol underperforms its random-policy baseline. Once the RLACO
initialization is available, Tabucol also shows the strongest overall
downstream refinement-time profile. Taken together, these results show that a
policy trained in one search framework can remain useful in another without
retraining, and that transferring its preferences through a population-based
construction process can convert that reusable guidance into improved final
graph-coloring performance and achieve the state-of-the-art results.

Future work should examine the transfer mechanism more systematically through
ablations of the reward design, learned policy, and soft-injection strength, in
order to identify which components contribute most strongly to transferable
guidance and favorable downstream initialization. The evaluation should also
be extended to larger and more diverse benchmark sets, alternative transfer
frameworks, and different downstream local-search solvers to determine how
broadly the observed transfer behavior generalizes. A further promising
direction is to reconsider the role of RL inside Tabucol itself. Rather than
using the policy only to propose recoloring actions, RL could act as a
trajectory-level monitor that observes the progress of Tabucol and dynamically
adjusts search parameters such as tabu tenure, diversification strength,
restart timing, or other search-control decisions. 